\documentclass[conference]{IEEEtran}
\usepackage{amsmath, amssymb}
\usepackage{graphicx}
\usepackage{booktabs}
\usepackage{cite}

\title{Multi-Robot Infield Defense: Formation, Interception Bidding, and Delivery Modeling in Simulation}

\author{
\IEEEauthorblockN{Aman Khatri, Dena Shink, Jason Ballinger}
\IEEEauthorblockA{Indiana University Bloomington}
}

\begin{document}
\maketitle

\begin{abstract}
We present a lightweight, end-to-end simulation of a baseball-inspired multi-robot infield defense task. A team of $N$ robots maintains a decentralized ring-graph formation, detects a rolling ground ball modeled by constant planar acceleration, assigns a primary fielder via distributed bidding, and then delivers the ball to a target base. Our focus is architectural clarity (what is computed locally vs. globally in a simulator harness) and a repeatable evaluation protocol. We implement the system in Python (NumPy + Matplotlib) and report repeated-trial results across $N$, $v_{\max}$, delivery distance, and handoff delay.
\end{abstract}

\section{Introduction}
Multi-robot systems (MRS) must often allocate tasks online while maintaining team structure. Inspired by baseball infield defense, we consider a simplified scenario where multiple robots must (i) maintain a defensive formation, (ii) intercept a moving ground ball, and (iii) record a force out against a runner advancing around the bases.

\textbf{Contribution.} Our contribution is a complete, runnable simulation pipeline that integrates:
\begin{itemize}
    \item decentralized formation control on a ring communication graph,
    \item per-robot interception bidding (including a moving-ball predicted bid),
    \item a runner-aware post-capture force-play model (direct vs. relay),
    \item a headless repeated-trial evaluation runner that generates a CSV for tables/plots.
\end{itemize}
We do not claim a new task-allocation primitive; rather, we make the architecture and evaluation protocol explicit and measurable.

\subsection{Related Work}
Market/auction-based task allocation is a standard approach for multi-robot assignment and provides a conceptual basis for bidding on task costs~\cite{gerkey2004}. Formation and consensus control using local neighbor interactions is also well established and can be described using graph Laplacians and agreement dynamics~\cite{mesbahi2010,olfati2006}. Our work combines these ideas in a deliberately simple setting to prioritize end-to-end completeness and quantitative evaluation in a controlled simulator.

\section{System Model}
Robots are holonomic point agents in 2D with positions \(p_i \in \mathbb{R}^2\) and bounded speed \(\|v_i\|\le v_{\max}\). The ball position is \(p_b(t)\in\mathbb{R}^2\). The objective is to minimize total completion time
\[
T_{\mathrm{total}} = T_{\mathrm{intercept}} + T_{\mathrm{delivery}}.
\]

\section{Methods}
\subsection{Ball trajectory model}
We approximate ground-ball motion with constant planar acceleration:
\begin{equation}
p_b(t) = p_0 + v_0 t + \frac{1}{2} a t^2, \qquad t \ge 0.
\end{equation}

\subsection{Formation control (local neighbor interactions)}
Let \(\mathcal{N}_i\) be robot \(i\)'s neighbors on a ring graph. We implement a Laplacian-style decentralized control law
\begin{equation}
u_i = -k \sum_{j \in \mathcal{N}_i} (p_i - p_j),
\label{eq:formation-law}
\end{equation}
using only neighbor positions. In the simulator, commands are clipped to respect \(v_{\max}\).

\subsection{Interception bidding and primary selection}
\textbf{Fixed-point bid.} At time \(t\), a simple travel-time proxy bid is
\begin{equation}
\widehat{T}_i(t) = \frac{\|p_i - p_b(t)\|}{v_{\max}}.
\label{eq:Ti-hat}
\end{equation}
\noindent This is easy to compute locally, but it treats the ball as a fixed target at time \(t\).

\textbf{Predicted moving-ball bid (implemented).} To account for ball motion, we implement a horizon-limited earliest feasible intercept time:
\begin{equation}
\tau_i^\ast = \min \left\{ \tau \in [0,H] : \|p_i - p_b(t+\tau)\| \le v_{\max}\tau \right\}.
\label{eq:pred-bid}
\end{equation}
Each robot can compute \(\tau_i^\ast\) given its own \(p_i\) and broadcast ball parameters \((p_0,v_0,a,t_0)\). In our implementation, \(\tau_i^\ast\) is computed by per-robot sampling in \([0,H]\) and a short binary search inside the first feasible bracket.

\noindent We then apply a lightweight baseball-style ownership bias: the infielder whose home position is closest to the ball's forward trajectory ray is treated as the trajectory owner, and non-owners incur a small score penalty unless they beat the owner's intercept time by a fixed margin. This preserves the moving-ball prediction while reducing unrealistic cross-field cutoffs on balls that stay in one fielder's lane.

\noindent The primary fielder is selected by the minimum biased fielding score, with a non-owner override only when its raw intercept time is sufficiently smaller than the owner's.
\begin{equation}
s_i = \tau_i^\ast + \lambda \,\mathbf{1}[i \neq i_{\mathrm{owner}}],
\end{equation}
\begin{equation}
i^\ast =
\begin{cases}
\arg\min_{i \neq i_{\mathrm{owner}}} \tau_i^\ast, & \text{if } \tau_i^\ast + \delta < \tau_{i_{\mathrm{owner}}}^\ast \text{ for some non-owner,}\\
i_{\mathrm{owner}}, & \text{otherwise when the owner is feasible,}\\
\arg\min_i s_i, & \text{otherwise.}
\end{cases}
\end{equation}
\noindent The selected primary then runs toward its predicted intercept point \(p_b(t+\tau_{i^\ast}^\ast)\) rather than the ball's instantaneous position.

\subsection{Delivery model: direct vs. relay}
After interception, we compare direct carry to a one-hop relay using straight-line legs at speed \(v_{\max}\). A batter-runner starts at home and advances around the basepath at the same speed bound. At each moment, the defense targets the runner's current force base \(p_{\mathrm{base}}\), and an out is recorded only if the ball carrier reaches that base before the runner. Let \(\tau_{\mathrm{hp}}\ge 0\) be a fixed handoff delay:
\begin{equation}
C_{\mathrm{direct}} = \frac{\|p_{i^\ast} - p_{\mathrm{base}}\|}{v_{\max}},
\end{equation}
\begin{equation}
C_{\mathrm{relay}}(j) = \frac{\|p_{i^\ast} - p_j\|}{v_{\max}} + \frac{\|p_j - p_{\mathrm{base}}\|}{v_{\max}} + \tau_{\mathrm{hp}}.
\end{equation}

\noindent\textbf{Limitation (triangle inequality).} Under this equal-speed running model, relays are not expected to strictly outperform direct carry (especially when \(\tau_{\mathrm{hp}}=0\)), because \(\|p_{i^\ast}-p_{\mathrm{base}}\|\le \|p_{i^\ast}-p_j\|+\|p_j-p_{\mathrm{base}}\|\). A more realistic baseball analogy would model a faster throw/pass speed than running; we discuss this in Section~\ref{sec:discussion}.

\section{Architecture: distributed rules in a centralized simulator}
Our \emph{implementation} is a centralized simulator harness that advances time, applies Euler integration, and renders plots. The \emph{decision rules} are written in a per-robot form that can be executed distributively with message passing.

\textbf{Explicit per-robot computation hooks (code).} The per-robot bid computations are implemented as:
\begin{itemize}
    \item \texttt{coordination.interception\_time\_estimate} (Eq.~\ref{eq:Ti-hat}),
    \item \texttt{coordination.interception\_time\_to\_moving\_ball} (Eq.~\ref{eq:pred-bid}),
    \item \texttt{coordination.plan\_fielding\_assignment} (trajectory-aware owner-biased assignment).
\end{itemize}
The simulator executes these computations centrally for convenience, but their inputs/outputs are per-agent and consistent with a distributed message-passing interpretation.

\section{Experimental Design}
\subsection{Protocol}
We use a headless evaluation runner (\texttt{python3 -m infield\_defense.evaluate}) that runs \(K=50\) randomized trials per setting and writes \texttt{results/eval.csv}. Ball launches are randomized within bounded ranges for initial position/velocity; each episode ends on a recorded out, a runner score, or a time limit.

\subsection{Baselines}
We compare three methods:
\begin{itemize}
    \item \textbf{nearest\_direct}: choose the nearest robot to the ball at detection, then deliver directly,
    \item \textbf{ours\_direct}: choose the primary by predicted moving-ball bid, then deliver directly,
    \item \textbf{ours\_relay}: choose the primary by predicted moving-ball bid, then choose direct vs. relay by the delivery cost model.
\end{itemize}

\subsection{Metrics}
We report success rate and total time \(T_{\mathrm{total}}\). We also log whether the relay option was selected (\texttt{relay\_used}).

\section{Results}
Table~\ref{tab:results-main} shows representative results for \(N=6\), \(v_{\max}=6\). ``base\_scale'' multiplies the nominal first-base location to study longer deliveries. ``Relay used'' reports the fraction of trials where the relay option was selected by the cost model.

\begin{table}[htbp]
\centering
\caption{Representative repeated-trial results (\(K=50\)) for \(N=6\), \(v_{\max}=6\). Metrics: success rate, relay-used rate, and \(T_{\mathrm{total}}\) (mean\(\pm\)std).}
\label{tab:results-main}
\begin{tabular}{@{}lccc@{}}
\toprule
Method & Success (\%) & Relay used (\%) & \(T_{\mathrm{total}}\) (s) \\
\midrule
\multicolumn{4}{@{}l@{}}{\textit{base\_scale = 1.0, \(\tau_{\mathrm{hp}}=0.0\)}} \\
nearest\_direct & 88.0 & 0.0 & 10.934 $\pm$ 0.246 \\
ours\_direct & 88.0 & 0.0 & 10.597 $\pm$ 0.432 \\
ours\_relay & 88.0 & 0.0 & 10.597 $\pm$ 0.432 \\
\midrule
\multicolumn{4}{@{}l@{}}{\textit{base\_scale = 2.0, \(\tau_{\mathrm{hp}}=0.0\)}} \\
nearest\_direct & 80.0 & 0.0 & 16.741 $\pm$ 0.259 \\
ours\_direct & 84.0 & 0.0 & 16.436 $\pm$ 0.529 \\
ours\_relay & 84.0 & 0.0 & 16.436 $\pm$ 0.529 \\
\midrule
\multicolumn{4}{@{}l@{}}{\textit{base\_scale = 2.0, \(\tau_{\mathrm{hp}}=0.5\)}} \\
nearest\_direct & 88.0 & 0.0 & 16.747 $\pm$ 0.236 \\
ours\_direct & 92.0 & 0.0 & 16.552 $\pm$ 0.472 \\
ours\_relay & 92.0 & 0.0 & 16.552 $\pm$ 0.472 \\
\bottomrule
\end{tabular}
\end{table}

\noindent In these representative regimes, predicted-bid selection improves \(T_{\mathrm{total}}\) versus the nearest-distance baseline, supporting the claim that bidding on an intercept-time proxy can outperform a purely geometric nearest rule. As expected under the equal-speed running delivery model, relay is not selected in these sweeps.

\section{Discussion and Limitations}
\label{sec:discussion}
Our simulation simplifies perception, dynamics, and ball handling. The most important modeling limitation for relay evaluation is that delivery legs are modeled with the same running speed. A more realistic relay would model a faster throw/pass speed \(v_{\mathrm{throw}}>v_{\max}\) plus a handoff time; this extension would create regimes where relay delivery outperforms direct carry, especially for longer base distances.

\section{Conclusion}
We implemented an end-to-end infield-defense-inspired MRS simulation with decentralized formation control, per-robot interception bidding (including a moving-ball predicted bid), and a headless evaluation harness that produces quantitative baseline comparisons. The current force-play model is transparent and measurable, and the next step for relay realism is incorporating a faster pass/throw model.

\begin{thebibliography}{00}
\bibitem{mesbahi2010} M. Mesbahi and M. Egerstedt, \textit{Graph Theoretic Methods in Multiagent Networks}, Princeton University Press, 2010.
\bibitem{gerkey2004} B. P. Gerkey and M. J. Mataric, ``A formal analysis and taxonomy of task allocation in multi-robot systems,'' \textit{International Journal of Robotics Research}, 2004.
\bibitem{olfati2006} R. Olfati-Saber, ``Flocking for multi-agent dynamic systems,'' \textit{IEEE Transactions on Automatic Control}, 2006.
\end{thebibliography}

\end{document}
